\documentclass[11pt]{article}

\usepackage{amsmath,amssymb}
\usepackage{xurl}
\usepackage[hidelinks]{hyperref}
\setlength{\emergencystretch}{2em}

% Shared commands for notes in docs/paper-gaps/.
% Notation follows the LDT paper (references/ldt-paper/).

\newcommand{\Lean}{\textsc{Lean}}
\newcommand{\leanid}[1]{\textsf{#1}}
\newcommand{\ghissue}[1]{\href{https://github.com/LionSR/MIPStarRE/issues/#1}{GitHub issue \##1}}
\newcommand{\ghpr}[1]{\href{https://github.com/LionSR/MIPStarRE/pull/#1}{GitHub PR \##1}}

% --- Paper-compatible notation ---
\newcommand{\F}{\mathbb{F}}
\newcommand{\N}{\mathbb{N}}
\newcommand{\C}{\mathbb{C}}
\newcommand{\tr}{\operatorname{tr}}
\newcommand{\ot}{\otimes}
\newcommand{\E}{\mathop{\bf E\/}}
\newcommand{\bx}{\mathbf{x}}
\newcommand{\by}{\mathbf{y}}
\newcommand{\bu}{\mathbf{u}}
\newcommand{\bv}{\mathbf{v}}

% Approximation relation (paper uses \approx_\eps and \simeq_\zeta)
\newcommand{\approxeps}{\approx_{\varepsilon}}
\newcommand{\simgeq}{\simeq_{\zeta}}

% Polynomial spaces (paper uses \polyfunc, \polysub, \polymeas)
\newcommand{\polyfunc}[3]{\operatorname{Poly}(#1,#2,#3)}
\newcommand{\polysub}[3]{\operatorname{PolySub}(#1,#2,#3)}
\newcommand{\polymeas}[3]{\operatorname{PolyMeas}(#1,#2,#3)}


\title{Review: How the LDT Paper Uses Singular Value Decomposition}
\author{}
\date{2026-05-05}

\begin{document}
\maketitle

\section{Where the SVD enters}

We examine the use of singular value decomposition in
\cite[Section~``Orthonormalization'']{Ji2020LowIndividualDegree}.%
\footnote{The SVD passage is in
\path{references/ldt-paper/orthonormalization.tex}, lines~774--845.  The
identities derived from it are in lines~862--896 and 984--1064, and the final
$P$-versus-$Q$ estimate using those identities is in lines~1089--1192.}
The literal SVD step appears only in the proof of the orthonormalization main
lemma.  Earlier parts of the proof use spectral truncation of the effects
$A_a$ and an orthonormal basis of the ranges of projectors; those are separate
finite-dimensional spectral arguments, not the SVD step reviewed here.

After rounding and rank reduction, the paper has a family of projectors
$Q=\{Q_a\}$ satisfying $\sum_a \operatorname{rank}(Q_a)\le d$.  For each
outcome $a$, it chooses an orthonormal basis
$\ket{v_{a,1}},\ldots,\ket{v_{a,m_a}}$ for the range of $Q_a$ and writes
\[
  Q_a = \sum_{i=1}^{m_a} \ket{v_{a,i}}\bra{v_{a,i}}.
\]
Let $m=\sum_a m_a$.  The rank-reduction bound gives $m\le d$.  On an auxiliary
space with orthonormal basis vectors $\ket{a,i}$, the paper defines
\[
  X_a = \sum_{i=1}^{m_a} \ket{a,i}\bra{v_{a,i}},
  \qquad
  X = \sum_a X_a,
  \tag{\texttt{orthonormalization.tex:784--790}}
\]
and a projective measurement $T=\{T_a\}$ by
\[
  T_a = \sum_{i=1}^{m_a}\ket{a,i}\bra{a,i}.
\]
The basic identities are
\[
  X_a = T_a X,
  \qquad
  Q_a = X_a^\dagger X_a = X^\dagger T_a X = X_a^\dagger X.
  \tag{\texttt{lem:xa-t}, \texttt{lem:qa-restated}}
\]
The SVD is then applied to this rectangular $m\times d$ matrix $X$:
\[
  X = U\Sigma_{m\times d}V^\dagger,
  \tag{\texttt{orthonormalization.tex:837--845}}
\]
where $U$ is $m\times m$ unitary, $V$ is $d\times d$ unitary, and
$\Sigma_{m\times d}$ is diagonal with nonnegative real entries.

\section{What the SVD is used to prove}

The SVD is not used as a black-box approximation theorem.  It is used to build a
single replacement matrix by changing the singular values of $X$ to $1$.  The
paper defines
\[
  \widehat X = U I_{m\times d} V^\dagger,
  \qquad
  \widehat X_a = T_a\widehat X,
  \qquad
  P_a = \widehat X_a^\dagger \widehat X_a.
  \tag{\texttt{orthonormalization.tex:954--963}}
\]
Because $m\le d$, the rectangular identity satisfies
$I_{m\times d}I_{d\times m}=I_{m\times m}$, so
\[
  \widehat X\widehat X^\dagger = I_{m\times m}.
  \tag{\texttt{lem:X-hat-squared}}
\]
This is the algebraic reason $P=\{P_a\}$ becomes a projective submeasurement:
\[
  P_aP_b
  = \widehat X^\dagger T_a\widehat X\widehat X^\dagger T_b\widehat X
  = \widehat X^\dagger T_aT_b\widehat X
  = \mathbf 1_{a=b}P_a.
  \tag{\texttt{lem:P-projectivity}}
\]

The second SVD-derived identity is the square-root comparison
\[
  X^\dagger \widehat X = \sqrt Q,
  \qquad Q=\sum_a Q_a.
  \tag{\texttt{lem:X-times-X-hat}}
\]
Indeed, multiplying the SVD formulas gives
$X^\dagger\widehat X=V\Sigma_{d\times d}V^\dagger$, while
$X^\dagger X=Q=V\Sigma_{d\times d}^2V^\dagger$.  Since the middle factor is
positive, it is the positive square root of $Q$.

The third use is the squared-difference estimate
\[
  (X-\widehat X)(X-\widehat X)^\dagger
  \le (XX^\dagger-I_{m\times m})^2.
  \tag{\texttt{lem:squared-difference}}
\]
In SVD coordinates this says
$(\Sigma-I)^2\le (\Sigma^2-I)^2$ on the nonnegative diagonal.  The estimate is
then inserted into a Cauchy--Schwarz comparison which replaces an occurrence of
$X$ by $\widehat X$ in the $Q_aP_a$ overlap.  Together with the almost
projectivity bound for $Q$, this proves
\[
  Q_a\ot I \approx_{30\zeta^{1/4}} P_a\ot I.
  \tag{\texttt{lem:P-Q-approx}}
\]
Combining this with the rank-reduction estimate
$A_a\ot I\approx_{12\sqrt\zeta} Q_a\ot I$ gives the orthonormalization main
lemma's final bound
\[
  A_a\ot I \approx_{84\zeta^{1/4}} P_a\ot I.
\]
Thus the paper uses SVD as a polar-repair step: the right Gram operator
$X^\dagger X=Q$ is kept, while the left singular directions are completed to a
coisometry $\widehat X$ whose row Gram is the identity.

\section{The blueprint comparison}

The informal blueprint follows this organization.  It introduces the matrix
$X$ in \path{def:matrix-decomposition-Q}, records the SVD step in
\path{def:svd-of-X}, and defines $P$ from
$\widehat X=UI_{m\times d}V^\dagger$ in \path{def:projective-P}.%
\footnote{See \path{blueprint/src/chapter/ch04_projective.tex}, lines~704--785.}
It then records the SVD-derived identities separately: $X^\dagger X=Q$ in
\path{lem:X-squared}, $\widehat X\widehat X^\dagger=I$ in
\path{lem:X-hat-squared}, $X^\dagger\widehat X=\sqrt Q$ in
\path{lem:X-times-X-hat}, the squared-difference estimate in
\path{lem:squared-difference}, and the final $P$-versus-$Q$ estimate in
\path{lem:P-Q-approx}.%
\footnote{These blueprint nodes are in
\path{blueprint/src/chapter/ch04_projective.tex}, lines~833--868 and
933--1153.}

The blueprint is more explicit than the paper about the formal interface.  Its
SVD node says that the formalization has two related layers: a reduced data
package storing exactly the $X$, $\widehat X$, and $P$ identities needed
downstream, and rectangular-SVD producers which derive those identities when
explicit matrices $U,V,\Sigma,I_{m\times d}$ are supplied.  This comparison is
faithful to the paper: the proof only needs the coisometry identity for
$\widehat X$, the mixed square-root identity, and the resulting estimates; it
never uses the SVD factors for any other purpose.

\section{The current formal statement}

The Lean development mirrors the blueprint separation.  The type named after the
paper's SVD stage is a reduced package of downstream consequences rather than a
record of all SVD factors.%
\footnote{The abbreviation \leanid{svdOfX} is defined as
\leanid{QXPLayerData} in
\path{MIPStarRE/LDT/MakingMeasurementsProjective/QXPLayer/Core.lean}, lines
184--191.}
A value of \leanid{QXPLayerData} stores the $Q$-layer, the matrices $X$ and
$\widehat X$, the restatement $Q_a=X^\dagger T_aX$, projectivity of the $Q_a$,
and the two primitive identities
\[
  \widehat X\widehat X^\dagger=I,
  \qquad
  X^\dagger\widehat X=\sqrt Q.
\]
It also proves internally that $X^\dagger X=Q$ from the $Q_a=X^\dagger T_aX$
restatement and the fact that $T$ is a measurement.%
\footnote{See \leanid{QXPLayerData.ofQLayerAndSvdIdentities} in
\path{MIPStarRE/LDT/MakingMeasurementsProjective/QXPLayer/Core.lean}, lines
202--269.}

The formal development also records the paper-facing rectangular-SVD algebra.
If matrices $U,V,S$ and a rectangular identity factor $I_{m\times d}$ satisfy
the unitary laws and $X=USV^\dagger$, then Lean proves the coisometry identity
for $UI_{m\times d}V^\dagger$ and the raw mixed-product formula.  A separate
positive-square-root uniqueness lemma turns the mixed product into
$\sqrt Q$ when the middle factor is positive and has square $Q$.%
\footnote{The relevant declarations are
\leanid{rectangularSvd_xHat_coisometry},
\leanid{rectangularSvd_xHat_mixed_raw},
\leanid{rectangularSvd_middle_eq_sqrt_of_square}, and
\leanid{exists_xHat_of_rectangularSvd} in
\path{MIPStarRE/LDT/MakingMeasurementsProjective/QXPLayerIdentities.lean},
lines~1040--1178.}
Those statements are exactly the matrix calculations that the paper performs
with $X=U\Sigma V^\dagger$ and
$\widehat X=UI_{m\times d}V^\dagger$.

Recent formalization work goes one step further and constructs the same
$\widehat X$ data through a positive-Gram, or rectangular polar-decomposition,
route.  From $X^\dagger X=Q$, positivity of $Q$, and the dimension inequality
$m\le d$, Lean constructs a rectangular coisometry $\widehat X$ such that
\[
  \widehat X\widehat X^\dagger=I,
  \qquad
  X^\dagger\widehat X=\sqrt Q.
\]
This is the same mathematical content that the paper obtains from an SVD, but
it is packaged in terms of spectral rows of the positive Gram operator rather
than a general rectangular-SVD existence theorem.%
\footnote{See \leanid{exists_xHat_of_positive_gram_spectrum},
\leanid{exists_xHat_of_sigmaFinRangeEmbedding_positiveGram}, and
\leanid{pQApprox_ofRankReductionSigmaRangePositiveGram} in
\path{MIPStarRE/LDT/MakingMeasurementsProjective/QXPLayerIdentities.lean}.  The
project-local dependency was tracked in \ghissue{1228}, following the older
SVD-interface work in \ghissue{1117} and \ghissue{652}.}
The downstream Lean estimates \leanid{xHatSquared}, \leanid{xTimesXHat},
\leanid{squaredDifference}, and \leanid{pQApprox} consume only these primitive
identities, just as the paper's later proof does.

\section{Wrappers, local algebra, and Mathlib inputs}

It is useful to separate packaging from mathematical content in this layer.  The
names \leanid{svdOfX}, \leanid{QXPLayerData}, and the constructors
\leanid{QXPLayerData.ofQLayerAndSvdIdentities},
\leanid{QXPLayerData.ofRankReductionAndSvdIdentities}, and their existential
variants are mostly wrappers.  They keep the paper's $Q\to X\to\widehat X\to P$
objects visible and discharge bookkeeping such as $X^\dagger X=Q$ from
$Q_a=X^\dagger T_aX$ and $\sum_aT_a=I$.  They do not themselves prove a general
SVD theorem.

The substantive project-local algebra is in two places.  First, the
rectangular-SVD adapters prove that, once matrices $U,V,S,I_{m\times d}$ with
the paper's SVD laws are supplied, the choice
$\widehat X=UI_{m\times d}V^\dagger$ has the coisometry and mixed-square-root
properties needed downstream.  Second, the positive-Gram construction proves the
same properties without asking for an explicit SVD record: it constructs the
rows of $\widehat X$ from the positive spectral subspace of $Q=X^\dagger X$ and
then completes them to a rectangular coisometry.

The positive-Gram construction is where Mathlib is used most heavily.  The
formal proof uses Mathlib's finite-dimensional inner-product and orthonormal
basis API for the singular-vector normalization calculation, the matrix-adjoint
interface relating $X^\dagger X$ to an adjoint composition, the Hermitian
spectral theorem through eigenvalues and eigenvector bases of a Hermitian
matrix, and the continuous functional calculus square root with its uniqueness
property.%
\footnote{Representative Mathlib-facing declarations used in this layer include
\leanid{Orthonormal}, \leanid{LinearMap.adjoint_inner_right},
\leanid{Matrix.toEuclideanLin_conjTranspose_mul_self},
\leanid{Matrix.IsHermitian.eigenvectorBasis},
\leanid{Matrix.IsHermitian.eigenvalues}, \leanid{CFC.sqrt}, and
\leanid{CFC.sqrt_unique}.  The local coisometry existence helper
\leanid{Matrix.exists_mul_conjTranspose_eq_one_of_card_le} in
\path{MIPStarRE/Quantum/FiniteHilbert.lean} is built from Mathlib's
\leanid{stdOrthonormalBasis} and \leanid{LinearMap.isometryOfOrthonormal}.}
What is not currently supplied directly by Mathlib is a rectangular complex SVD
or polar-decomposition theorem with exactly the output shape needed here.  The
project therefore proves the necessary finite-dimensional positive-Gram route
locally.  Several of these generic linear-algebra lemmas may be suitable for
future Mathlib upstreaming, but the paper-gap verdict does not depend on that
upstreaming.%
\footnote{The possible upstreaming of generic project-local matrix lemmas is
tracked separately in \ghissue{1250}.}

\section{Conclusion}

The verdict is that there is no mathematical discrepancy in the paper's use of
SVD.  The SVD is a local finite-dimensional tool in the orthonormalization proof:
it produces $\widehat X$ by replacing the singular values of the rank-reduced
matrix $X$ with $1$, and it supplies the coisometry, square-root, and
squared-difference identities needed to prove that the projective family
$P_a=\widehat X^\dagger T_a\widehat X$ is close to $Q_a$.

The blueprint and Lean formalization preserve this mechanism.  Lean does not
need to store a full SVD record once the two primitive $\widehat X$ identities
are available, and it now has a positive-Gram construction producing those
identities in the sigma-range setting.  Remaining orthonormalization work, such
as supplying the broader spectral-truncation or locality-preserving repair
inputs, is outside this SVD step and should not be classified as a new paper
SVD gap.

\bibliographystyle{alpha}
\bibliography{references}

\end{document}
